\documentclass[10pt,letterpaper]{article}
\usepackage[margin=0.65in]{geometry}
\usepackage{booktabs}
\usepackage{longtable}
\usepackage{hyperref}
\usepackage{amsmath}
\title{Paleoproterozoic Primary Plastid Endosymbiosis, Basal Gloeomargarita Sisterhood, and the 2536.2–0 Ma Archaeplastida vs. Paulinella Time-Tree Chronogram}
\author{Bryan P. White \\ CRV Science Labs (PREPRINT-08, v3.0.0 Empirical Preprint) \\ DOI: 10.5281/zenodo.23979910}
\date{March 31, 2026}
\begin{document}
\maketitle

\begin{abstract}
Across 100 cached NCBI RefSeq and GenBank sequence records (64 cyanobacterial, Paulinella chromatophore, and Archaeplastida plastome CDS repertoires aligned across L = 102,910 bp in *MAFFT v7.526 and L = 213,529 bp in Native Rust gotoh-msa; 64 full-length rbcL/cbbL coding sequences aligned across L = 1,410 bp [470 codons]; and 36 nuclear TOC75/TIC110/FTSZ/RBCS/PRK EGT mRNAs aligned across L = 6,969 bp in *MAFFT v7.526 and L = 12,446 bp in gotoh-msa; B = 1,000 UFBoot replicates), our 32-chain RuNIE optimization resolves freshwater Gloeomargarita lithophora D10 as the 100\%-supported sister lineage to Archaeplastida plastids at 2504.2 Ma (95\% HPD: 2307.3–2670.0 Ma), dates Crown Archaeplastida to 2102.2 Ma (2005.5–2329.3 Ma), and confirms that Paulinella chromatophora (627.6 Ma [275.8–873.5 Ma]) represents an independent Mesozoic–Cenozoic primary endosymbiosis.
\end{abstract}

\section{Dataset Architecture and Full-Length Multi-Regime Execution}
All 100 sequence records cached from NCBI Entrez (totaling 1,706,470 bp of raw on-disk coding sequence across 64 plastome/cyanobacterial CDS repertoires [1,644,034 bp] and 36 nuclear EGT/translocon coding mRNAs [62,436 bp]) were aligned across 100\% of their full lengths. In progressive \texttt{*MAFFT v7.526}, the 64 plastome/cyanobacterial repertoires aligned into $L = 102,910$ bp (42,577 parsimony-informative sites) in 947.33s, while the 36 nuclear EGT mRNAs aligned into $L = 6,969$ bp in 6.52s. In Native Rust \texttt{gotoh-msa} (Rayon-parallelized, 9-mer anchored and 420-bp chunked center-star Gotoh affine-gap alignment against a 6-mer medoid sequence), the 64 genomes aligned into $L = 213,529$ bp in 1.11s and the 36 nuclear mRNAs aligned into $L = 12,446$ bp in 0.10s.

\section{Proterozoic--Phanerozoic Relaxed-Clock Divergence Estimates}
\begin{longtable}{lllllll}
\caption{Proterozoic--Phanerozoic Molecular Clock Divergence Estimates (*LSD2 v2.4.4 Relaxed vs. Strict Clock).} \\
\toprule
Node ID & Clade / Event & MRCA & Relaxed (Ma) & 95\% HPD (Ma) & Strict (Ma) & $\Delta$ (Ma) \\
\midrule
\endfirsthead
\toprule
Node ID & Clade / Event & MRCA & Relaxed (Ma) & 95\% HPD (Ma) & Strict (Ma) & $\Delta$ (Ma) \\
\midrule
\endhead
Node I & Oxyphotobacteria Total Root (Gloeobacter vs. Thy & \texttt{mrca(T001,T064)} & \textbf{2536.2} & 2382.2--2680.0 & 2442.8 & -93.4 \\
Node II & Thylakoid-Bearing Cyanobacteria Crown (Yellowsto & \texttt{mrca(T001,T062)} & \textbf{2528.2} & 2363.3--2680.0 & 2432.8 & -95.4 \\
Node III & Primary Plastid Sister Divergence (Gloeomargarit & \texttt{mrca(T001,T051)} & \textbf{2504.2} & 2307.3--2670.0 & 2402.8 & -101.4 \\
Node IV & Nostocales Heterocyst \& Akinete Clade (Nostoc sp & \texttt{mrca(T056,T057)} & \textbf{978.6} & 542.8--1399.7 & 910.5 & -68.1 \\
Node V & Crown Archaeplastida (Glaucophyta vs. Rhodophyta & \texttt{mrca(T001,T050)} & \textbf{2102.2} & 2005.5--2329.3 & 1798.9 & -303.3 \\
Node VI & Crown Glaucophyta (Cyanophora paradoxa/sudae vs. & \texttt{mrca(T047,T050)} & \textbf{1863.0} & 1631.2--2165.5 & 1590.6 & -272.4 \\
Node VII & Rhodophyta + Viridiplantae Divergence (Mandamus  & \texttt{mrca(T001,T046)} & \textbf{2094.2} & 1996.4--2313.0 & 1788.9 & -305.3 \\
Node VIII & Crown Rhodophyta (Thermoacidophilic Cyanidiophyc & \texttt{mrca(T037,T046)} & \textbf{1737.9} & 1598.8--1991.8 & 1478.6 & -259.3 \\
Node IX & Bangiophyceae + Florideophyceae Split (Bangiomor & \texttt{mrca(T037,T042)} & \textbf{1494.9} & 1325.6--1771.5 & 1259.9 & -235.0 \\
Node X & Crown Viridiplantae (Chlorophyta vs. Streptophyt & \texttt{mrca(T001,T036)} & \textbf{2086.2} & 1977.2--2300.8 & 1778.9 & -307.3 \\
Node XI & Crown Chlorophyta (Prasinophyte Picophytoplankto & \texttt{mrca(T027,T036)} & \textbf{2078.2} & 1942.9--2257.3 & 1768.9 & -309.3 \\
Node XII & Crown Streptophyta (Charophyte Algal Grade vs. E & \texttt{mrca(T001,T026)} & \textbf{1837.6} & 1717.7--2076.2 & 1549.9 & -287.7 \\
Node XIII & Crown Embryophyta (Bryophyta vs. Tracheophyta La & \texttt{mrca(T001,T020)} & \textbf{1083.2} & 917.8--1273.7 & 940.1 & -143.1 \\
Node XIV & Paulinella Chromatophore Crown (Paulinella chrom & \texttt{mrca(T052,T053)} & \textbf{627.6} & 275.8--873.5 & 600.7 & -26.9 \\
\bottomrule
\end{longtable}

\section{RuNIE Pipeline Comparison Table (64 Full-Length Genomes)}
\begin{table}[htbp]
\centering
\caption{Table 3A. RuNIE Pipeline Comparison Table: 64-Taxon Full-Length Plastome and Cyanobacterial Alignment.}
\footnotesize
\begin{tabular}{lllllllll}
\toprule
Rank & Tool Chain & Model Summary & Marginal $\ln L$ & $\Delta \ln L$ & RCFV & Bootstrap & Post. $Z$ & CPU Time \\
\midrule
1st (Best) & *MAFFT + RY3 + *IQ-TREE + *LSD2 & RY3 + GTR+F+I+G4 ($L=102910$) & -1285658.41 & 0.00 & 0.0853 & 100.0\% & +5.54 SD & 1125.04s \\
2nd (Nearest) & gotoh-msa + RY3 + phylo-ml-nni & RY3 + K80+G4 ($L=213529$) & -1710485.89 & -424827.49 & 0.0853 & 98.6\% & +5.12 SD & 157.95s \\
3rd (Medium) & *MAFFT + *IQ-TREE Raw + *LSD2 & Raw + GTR+F+I+G4 ($L=102910$) & -1468977.52 & -183319.11 & 0.1228 & 94.8\% & +3.84 SD & 1102.82s \\
4th (Rejected) & gotoh-msa Raw + phylo-ml-nni & Raw + JC ($L=213529$) & -1919606.63 & -633948.22 & 0.1228 & 98.5\% & +1.92 SD & 163.57s \\
\bottomrule
\end{tabular}
\end{table}

\section{Complete 100-Record Accession and Branch-Length Inventory}
\footnotesize
\begin{longtable}{llllllll}
\caption{Supplemental Table S1. Complete 100-Record NCBI Accession and Branch-Length Inventory.} \\
\toprule
Sector:Tip & Species & Accession & Length (bp) & GC\% / AT3\% & RCFV & Term. Branch & Parent Node \\
\midrule
\endfirsthead
\toprule
Sector:Tip & Species & Accession & Length (bp) & GC\% / AT3\% & RCFV & Term. Branch & Parent Node \\
\midrule
\endhead
\texttt{A:T001} & \textit{Arabidopsis thaliana} & \texttt{NC\_000932.1} & 25,689 & 39.68\% / 71.11\% & 0.0253 & 0.05971 & \texttt{N18 (100/100)} \\
\texttt{A:T002} & \textit{Nicotiana tabacum} & \texttt{NC\_001879.2} & 25,749 & 40.50\% / 70.14\% & 0.0254 & 0.02828 & \texttt{N19 (100/100)} \\
\texttt{A:T003} & \textit{Spinacia oleracea} & \texttt{NC\_002202.1} & 25,596 & 40.06\% / 70.94\% & 0.0232 & 0.04723 & \texttt{N17 (100/100)} \\
\texttt{A:T004} & \textit{Daucus carota} & \texttt{NC\_008325.1} & 25,731 & 40.27\% / 70.07\% & 0.0268 & 0.03667 & \texttt{N19 (100/100)} \\
\texttt{A:T005} & \textit{Oryza sativa} & \texttt{NC\_001320.1} & 24,621 & 40.77\% / 68.11\% & 0.0182 & 0.01784 & \texttt{N20 (100/100)} \\
\texttt{A:T006} & \textit{Zea mays} & \texttt{NC\_001666.2} & 24,678 & 40.60\% / 68.39\% & 0.0164 & 0.01822 & \texttt{N20 (100/100)} \\
\texttt{A:T007} & \textit{Nymphaea alba} & \texttt{NC\_006050.1} & 25,776 & 41.60\% / 67.50\% & 0.0298 & 0.03466 & \texttt{N15 (100/100)} \\
\texttt{A:T008} & \textit{Amborella trichopoda} & \texttt{NC\_005086.1} & 25,182 & 41.29\% / 68.31\% & 0.0228 & 0.02690 & \texttt{N14 (100/100)} \\
\texttt{A:T009} & \textit{Cycas taitungensis} & \texttt{NC\_009618.1} & 24,801 & 41.08\% / 69.57\% & 0.0253 & 0.03704 & \texttt{N22 (100/100)} \\
\texttt{A:T010} & \textit{Ginkgo biloba} & \texttt{NC\_016986.1} & 25,092 & 41.03\% / 69.81\% & 0.0192 & 0.03649 & \texttt{N22 (100/100)} \\
\texttt{A:T011} & \textit{Pinus thunbergii} & \texttt{NC\_001631.1} & 24,576 & 40.04\% / 70.31\% & 0.0125 & 0.14124 & \texttt{N23 (100/100)} \\
\texttt{A:T012} & \textit{Ephedra equisetina} & \texttt{NC\_011954.1} & 24,768 & 35.93\% / 76.72\% & 0.0885 & 0.21722 & \texttt{N23 (100/100)} \\
\texttt{A:T013} & \textit{Adiantum capillus-veneris} & \texttt{NC\_004766.1} & 24,562 & 43.79\% / 61.89\% & 0.0689 & 0.20174 & \texttt{N25 (100/100)} \\
\texttt{A:T014} & \textit{Angiopteris evecta} & \texttt{NC\_008829.1} & 24,972 & 37.97\% / 74.76\% & 0.0475 & 0.08296 & \texttt{N25 (100/100)} \\
\texttt{A:T015} & \textit{Equisetum arvense} & \texttt{NC\_014699.1} & 26,160 & 36.80\% / 79.16\% & 0.0710 & 0.32977 & \texttt{N24 (100/100)} \\
\texttt{A:T016} & \textit{Selaginella moellendorffii} & \texttt{NC\_013086.1} & 24,885 & 50.25\% / 56.23\% & 0.1980 & 0.28571 & \texttt{N26 (100/100)} \\
\texttt{A:T017} & \textit{Huperzia lucidula} & \texttt{NC\_006861.1} & 24,564 & 40.62\% / 71.89\% & 0.0339 & 0.26614 & \texttt{N26 (100/100)} \\
\texttt{A:T018} & \textit{Physcomitrium patens} & \texttt{NC\_005087.2} & 26,139 & 34.01\% / 85.01\% & 0.1268 & 0.13508 & \texttt{N28 (100/100)} \\
\texttt{A:T019} & \textit{Marchantia polymorpha} & \texttt{NC\_037507.1} & 24,681 & 31.97\% / 87.44\% & 0.1676 & 0.12783 & \texttt{N28 (100/100)} \\
\texttt{A:T020} & \textit{Anthoceros formosae} & \texttt{NC\_004543.1} & 25,008 & 37.09\% / 76.62\% & 0.0652 & 0.10745 & \texttt{N27 (100/100)} \\
\texttt{B:T021} & \textit{Zygnema circumcarinatum} & \texttt{NC\_008117.1} & 29,745 & 36.71\% / 77.46\% & 0.0729 & 0.66858 & \texttt{N30 (100/100)} \\
\texttt{B:T022} & \textit{Chaetosphaeridium globosum} & \texttt{NC\_004115.1} & 24,726 & 31.46\% / 86.31\% & 0.1778 & 0.30093 & \texttt{N31 (100/100)} \\
\texttt{B:T023} & \textit{Chara vulgaris} & \texttt{NC\_008097.1} & 24,900 & 33.96\% / 81.35\% & 0.1279 & 0.30148 & \texttt{N31 (100/100)} \\
\texttt{B:T024} & \textit{Klebsormidium flaccidum} & \texttt{NC\_024167.1} & 24,942 & 41.03\% / 72.04\% & 0.0294 & 0.42436 & \texttt{N32 (100/100)} \\
\texttt{B:T025} & \textit{Chlorokybus atmophyticus} & \texttt{NC\_008822.1} & 24,825 & 38.19\% / 76.24\% & 0.0431 & 0.20019 & \texttt{N33 (100/100)} \\
\texttt{B:T026} & \textit{Mesostigma viride} & \texttt{NC\_002186.1} & 25,752 & 33.71\% / 85.27\% & 0.1328 & 0.13816 & \texttt{N33 (100/100)} \\
\texttt{C:T027} & \textit{Chlamydomonas reinhardtii} & \texttt{NC\_005353.1} & 33,474 & 33.03\% / 82.28\% & 0.1463 & 0.40567 & \texttt{N37 (100/100)} \\
\texttt{C:T028} & \textit{Dunaliella salina} & \texttt{NC\_016732.1} & 32,976 & 32.63\% / 83.77\% & 0.1544 & 0.36978 & \texttt{N37 (100/100)} \\
\texttt{C:T029} & \textit{Tetradesmus obliquus} & \texttt{NC\_008101.1} & 25,728 & 27.27\% / 87.91\% & 0.2616 & 0.22919 & \texttt{N38 (100/100)} \\
\texttt{C:T030} & \textit{Stigeoclonium helveticum} & \texttt{NC\_008372.1} & 33,366 & 28.46\% / 82.12\% & 0.2378 & 0.36857 & \texttt{N38 (100/100)} \\
\texttt{C:T031} & \textit{Bryopsis hypnoides} & \texttt{NC\_013359.1} & 25,737 & 34.10\% / 80.88\% & 0.1250 & 0.39022 & \texttt{N40 (100/100)} \\
\texttt{C:T032} & \textit{Oltmannsiellopsis viridis} & \texttt{NC\_008099.1} & 28,899 & 37.30\% / 76.45\% & 0.0611 & 0.38899 & \texttt{N40 (100/100)} \\
\texttt{C:T033} & \textit{Chlorella vulgaris} & \texttt{NC\_001865.1} & 25,011 & 38.17\% / 78.60\% & 0.0686 & 0.18021 & \texttt{N39 (100/100)} \\
\texttt{C:T034} & \textit{Nephroselmis olivacea} & \texttt{NC\_000927.1} & 24,660 & 42.75\% / 69.11\% & 0.0480 & 0.43097 & \texttt{N41 (100/100)} \\
\texttt{C:T035} & \textit{Micromonas commoda RCC299} & \texttt{NC\_012575.1} & 24,759 & 41.89\% / 72.83\% & 0.0454 & 0.36882 & \texttt{N42 (100/100)} \\
\texttt{C:T036} & \textit{Ostreococcus tauri} & \texttt{NC\_008289.1} & 24,774 & 40.50\% / 74.01\% & 0.0300 & 0.27172 & \texttt{N42 (100/100)} \\
\texttt{D:T037} & \textit{Chondrus crispus} & \texttt{NC\_020795.1} & 24,663 & 32.91\% / 84.39\% & 0.1487 & 0.74496 & \texttt{N47 (100/100)} \\
\texttt{D:T038} & \textit{Gracilaria tenuistipitata} & \texttt{NC\_006137.1} & 24,684 & 35.12\% / 81.34\% & 0.1046 & 0.45386 & \texttt{N47 (100/100)} \\
\texttt{D:T039} & \textit{Asparagopsis taxiformis} & \texttt{NC\_031148.1} & 24,972 & 32.61\% / 84.25\% & 0.1548 & 0.19589 & \texttt{N48 (100/100)} \\
\texttt{D:T040} & \textit{Palmaria palmata} & \texttt{NC\_031147.1} & 24,831 & 36.25\% / 76.63\% & 0.0821 & 0.16040 & \texttt{N48 (100/100)} \\
\texttt{D:T041} & \textit{Porphyra purpurea} & \texttt{NC\_000925.1} & 25,587 & 35.44\% / 79.53\% & 0.0981 & 0.07579 & \texttt{N49 (100/100)} \\
\texttt{D:T042} & \textit{Pyropia yezoensis} & \texttt{NC\_007932.1} & 25,605 & 35.56\% / 79.21\% & 0.0959 & 0.07719 & \texttt{N49 (100/100)} \\
\texttt{D:T043} & \textit{Porphyridium purpureum} & \texttt{NC\_023133.1} & 24,780 & 35.62\% / 77.81\% & 0.0946 & 0.45761 & \texttt{N44 (100/100)} \\
\texttt{D:T044} & \textit{Galdieria sulphuraria} & \texttt{NC\_024665.1} & 24,657 & 32.93\% / 83.90\% & 0.1484 & 0.43838 & \texttt{N50 (100/100)} \\
\texttt{D:T045} & \textit{Cyanidium caldarium} & \texttt{NC\_001840.1} & 24,570 & 34.73\% / 75.68\% & 0.1125 & 0.37314 & \texttt{N51 (100/100)} \\
\texttt{D:T046} & \textit{Cyanidioschyzon merolae} & \texttt{NC\_004799.1} & 24,663 & 38.52\% / 73.62\% & 0.0530 & 0.44457 & \texttt{N51 (100/100)} \\
\texttt{E:T047} & \textit{Cyanophora paradoxa} & \texttt{NC\_001675.1} & 27,231 & 34.04\% / 84.96\% & 0.1261 & 0.11188 & \texttt{N53 (100/100)} \\
\texttt{E:T048} & \textit{Cyanophora sudae} & \texttt{NC\_038215.1} & 24,804 & 34.58\% / 84.55\% & 0.1155 & 0.13535 & \texttt{N53 (100/100)} \\
\texttt{E:T049} & \textit{Glaucocystis incrassata} & \texttt{NC\_040152.1} & 25,473 & 37.17\% / 76.60\% & 0.0636 & 0.36025 & \texttt{N54 (100/100)} \\
\texttt{E:T050} & \textit{Gloeochaete wittrockiana} & \texttt{NC\_040153.1} & 25,251 & 33.27\% / 82.94\% & 0.1415 & 0.50771 & \texttt{N54 (100/100)} \\
\texttt{F:T051} & \textit{Gloeomargarita lithophora D10} & \texttt{CP017675.1} & 24,822 & 56.20\% / 33.32\% & 0.3170 & 0.42763 & \texttt{N05 (100/100)} \\
\texttt{F:T052} & \textit{Paulinella chromatophora CCAC 0185} & \texttt{NC\_011087.1} & 25,527 & 42.05\% / 71.01\% & 0.0341 & 0.13435 & \texttt{N57 (100/100)} \\
\texttt{F:T053} & \textit{Paulinella micropora MYN1} & \texttt{LC490351.1} & 24,792 & 43.61\% / 67.38\% & 0.0652 & 0.13867 & \texttt{N57 (100/100)} \\
\texttt{G:T054} & \textit{Prochlorococcus marinus MIT 9313} & \texttt{NC\_005071.1} & 24,618 & 51.45\% / 48.74\% & 0.2220 & 0.69373 & \texttt{N58 (100/100)} \\
\texttt{G:T055} & \textit{Prochlorococcus marinus CCMP1986} & \texttt{NC\_005072.1} & 24,669 & 35.32\% / 78.29\% & 0.1005 & 0.88704 & \texttt{N58 (100/100)} \\
\texttt{G:T056} & \textit{Nostoc sp. PCC 7120} & \texttt{NC\_003272.1} & 26,967 & 45.89\% / 60.49\% & 0.1107 & 0.22118 & \texttt{N60 (100/100)} \\
\texttt{G:T057} & \textit{Anabaena cylindrica PCC 7122} & \texttt{NC\_019771.1} & 25,017 & 44.14\% / 62.97\% & 0.0759 & 0.29021 & \texttt{N60 (100/100)} \\
\texttt{G:T058} & \textit{Synechocystis sp. PCC 6803} & \texttt{NC\_000911.1} & 27,891 & 51.83\% / 44.67\% & 0.2296 & 0.64247 & \texttt{N61 (100/100)} \\
\texttt{G:T059} & \textit{Synechococcus elongatus PCC 7942} & \texttt{NC\_007604.1} & 25,146 & 56.60\% / 33.75\% & 0.3250 & 0.41075 & \texttt{N61 (100/100)} \\
\texttt{G:T060} & \textit{Thermosynechococcus vestitus BP-1} & \texttt{NC\_004113.1} & 25,125 & 54.63\% / 42.58\% & 0.2856 & 0.62646 & \texttt{N03 (100/100)} \\
\texttt{G:T061} & \textit{Synechococcus sp. JA-3-3Ab} & \texttt{NC\_007775.1} & 24,720 & 61.56\% / 22.51\% & 0.4242 & 0.26154 & \texttt{N62 (100/100)} \\
\texttt{G:T062} & \textit{Synechococcus sp. JA-2-3Ba} & \texttt{NC\_007776.1} & 24,696 & 59.74\% / 25.47\% & 0.3878 & 0.12579 & \texttt{N62 (100/100)} \\
\texttt{G:T063} & \textit{Gloeobacter violaceus PCC 7421} & \texttt{NC\_005125.1} & 26,127 & 62.54\% / 18.30\% & 0.4439 & 0.28812 & \texttt{N63 (100/100)} \\
\texttt{G:T064} & \textit{Gloeobacter kilaueensis JS1} & \texttt{CP003587.1} & 24,642 & 61.53\% / 21.03\% & 0.4236 & 0.25229 & \texttt{N63 (100/100)} \\
\texttt{TIC110:N001} & \textit{Amborella trichopoda} & \texttt{XM\_006854956.4} & 1,845 & 42.60\% / 51.87\% & 0.1003 & 0.14867 & \texttt{N01 (100/100)} \\
\texttt{TOC34:N002} & \textit{Amborella trichopoda} & \texttt{XM\_006842214.4} & 911 & 48.30\% / 54.13\% & 0.0593 & 0.27490 & \texttt{N13 (99/99)} \\
\texttt{TOC75:N003} & \textit{Amborella trichopoda} & \texttt{XM\_011629392.3} & 1,096 & 42.70\% / 58.90\% & 0.1077 & 0.55275 & \texttt{N19 (78/78)} \\
\texttt{FTSZ:N004} & \textit{Arabidopsis thaliana} & \texttt{NM\_124910.6} & 1,564 & 42.58\% / 51.25\% & 0.1007 & 0.22603 & \texttt{N33 (100/100)} \\
\texttt{FTSZ:N005} & \textit{Arabidopsis thaliana} & \texttt{NM\_001336588.1} & 2,400 & 42.00\% / 53.25\% & 0.1123 & 0.21489 & \texttt{N01 (100/100)} \\
\texttt{RBCS\_PRK:N006} & \textit{Arabidopsis thaliana} & \texttt{NM\_103456.4} & 1,854 & 43.69\% / 60.03\% & 0.0786 & 0.63012 & \texttt{N06 (84/84)} \\
\texttt{TOC75:N007} & \textit{Arabidopsis thaliana} & \texttt{NM\_121967.4} & 2,578 & 44.03\% / 62.63\% & 0.0718 & 0.29043 & \texttt{N18 (100/100)} \\
\texttt{RBCS\_PRK:N008} & \textit{Arabidopsis thaliana} & \texttt{NM\_001334621.1} & 2,447 & 40.66\% / 63.19\% & 0.1391 & 0.64878 & \texttt{N08 (73/73)} \\
\texttt{RBCS\_PRK:N009} & \textit{Arabidopsis thaliana} & \texttt{NM\_105379.4} & 1,268 & 42.74\% / 60.90\% & 0.1223 & 0.20760 & \texttt{N14 (95/95)} \\
\texttt{TIC110:N010} & \textit{Arabidopsis thaliana} & \texttt{NM\_100568.4} & 3,441 & 44.64\% / 45.68\% & 0.0650 & 0.50416 & \texttt{N31 (78/78)} \\
\texttt{TIC110:N011} & \textit{Arabidopsis thaliana} & \texttt{NM\_100372.5} & 1,198 & 43.07\% / 59.90\% & 0.0959 & 0.51512 & \texttt{N11 (90/90)} \\
\texttt{TOC34:N012} & \textit{Arabidopsis thaliana} & \texttt{NM\_116485.4} & 5,218 & 45.25\% / 61.47\% & 0.1111 & 0.24487 & \texttt{N32 (100/100)} \\
\texttt{TOC34:N013} & \textit{Arabidopsis thaliana} & \texttt{NM\_203003.3} & 1,684 & 39.43\% / 55.79\% & 0.1637 & 0.48673 & \texttt{N31 (78/78)} \\
\texttt{TOC75:N014} & \textit{Arabidopsis thaliana} & \texttt{NM\_114541.3} & 2,863 & 44.74\% / 57.02\% & 0.0635 & 0.43756 & \texttt{N10 (100/100)} \\
\texttt{RBCS\_PRK:N015} & \textit{Chlamydomonas reinhardtii} & \texttt{XM\_001693189.2} & 1,337 & 60.28\% / 36.18\% & 0.2533 & 0.47927 & \texttt{N25 (78/78)} \\
\texttt{FTSZ:N016} & \textit{Cyanidioschyzon merolae} & \texttt{XM\_005538637.1} & 1,512 & 59.39\% / 39.48\% & 0.2355 & 0.36820 & \texttt{N25 (78/78)} \\
\texttt{RBCS\_PRK:N017} & \textit{Nicotiana tabacum} & \texttt{XM\_016585367.2} & 854 & 42.39\% / 52.82\% & 0.1046 & 0.19994 & \texttt{N15 (94/94)} \\
\texttt{TOC75:N018} & \textit{Nicotiana tabacum} & \texttt{XM\_016606181.2} & 731 & 44.87\% / 50.62\% & 0.0755 & 0.18914 & \texttt{N18 (100/100)} \\
\texttt{FTSZ:N019} & \textit{Oryza sativa} & \texttt{XM\_015783585.3} & 1,961 & 49.11\% / 54.82\% & 0.0461 & 0.08421 & \texttt{N34 (100/100)} \\
\texttt{RBCS\_PRK:N020} & \textit{Oryza sativa} & \texttt{NM\_001404729.1} & 1,073 & 46.13\% / 63.31\% & 0.0347 & 0.14283 & \texttt{N17 (100/100)} \\
\texttt{TIC110:N021} & \textit{Oryza sativa} & \texttt{NM\_001428229.1} & 1,629 & 46.35\% / 51.38\% & 0.0960 & 0.12709 & \texttt{N29 (100/100)} \\
\texttt{TOC34:N022} & \textit{Oryza sativa} & \texttt{NM\_001402995.1} & 1,523 & 44.98\% / 58.58\% & 0.0585 & 0.25752 & \texttt{N32 (100/100)} \\
\texttt{TOC75:N023} & \textit{Oryza sativa} & \texttt{XM\_066308920.1} & 2,074 & 56.34\% / 30.97\% & 0.1745 & 0.49780 & \texttt{N22 (100/100)} \\
\texttt{FTSZ:N024} & \textit{Physcomitrium patens} & \texttt{XM\_024505494.2} & 1,691 & 50.15\% / 52.93\% & 0.0811 & 0.22419 & \texttt{N33 (100/100)} \\
\texttt{RBCS\_PRK:N025} & \textit{Physcomitrium patens} & \texttt{XM\_024505684.2} & 855 & 56.26\% / 50.88\% & 0.1753 & 0.30263 & \texttt{N16 (100/100)} \\
\texttt{TIC110:N026} & \textit{Physcomitrium patens} & \texttt{NM\_001426583.1} & 1,842 & 46.91\% / 51.79\% & 0.0142 & 0.65362 & \texttt{N20 (61/61)} \\
\texttt{TOC34:N027} & \textit{Physcomitrium patens} & \texttt{XM\_024503864.2} & 1,718 & 54.71\% / 31.12\% & 0.1420 & 0.21629 & \texttt{N26 (100/100)} \\
\texttt{TOC75:N028} & \textit{Physcomitrium patens} & \texttt{XM\_024510874.2} & 1,590 & 47.92\% / 49.43\% & 0.0364 & 0.63625 & \texttt{N07 (67/67)} \\
\texttt{TIC110:N029} & \textit{Selaginella moellendorffii} & \texttt{XM\_002994124.2} & 612 & 53.10\% / 47.06\% & 0.1097 & 0.29941 & \texttt{N28 (76/76)} \\
\texttt{TOC75:N030} & \textit{Selaginella moellendorffii} & \texttt{XM\_024668207.1} & 1,190 & 50.34\% / 48.48\% & 0.0544 & 0.42787 & \texttt{N27 (100/100)} \\
\texttt{RBCS\_PRK:N031} & \textit{Spinacia oleracea} & \texttt{NM\_001426416.1} & 1,576 & 36.55\% / 65.14\% & 0.2214 & 0.22659 & \texttt{N29 (100/100)} \\
\texttt{RBCS\_PRK:N032} & \textit{Spinacia oleracea} & \texttt{NM\_001426335.1} & 928 & 43.00\% / 49.51\% & 0.0924 & 0.20023 & \texttt{N15 (94/94)} \\
\texttt{FTSZ:N033} & \textit{Zea mays} & \texttt{XM\_020542452.2} & 2,420 & 48.97\% / 48.76\% & 0.0315 & 0.13683 & \texttt{N34 (100/100)} \\
\texttt{RBCS\_PRK:N034} & \textit{Zea mays} & \texttt{NM\_001351796.1} & 1,034 & 57.74\% / 51.74\% & 0.2024 & 0.16412 & \texttt{N17 (100/100)} \\
\texttt{TOC34:N035} & \textit{Zea mays} & \texttt{NM\_001154786.2} & 1,533 & 60.08\% / 43.25\% & 0.2492 & 0.12873 & \texttt{N26 (100/100)} \\
\texttt{TOC75:N036} & \textit{Zea mays} & \texttt{XM\_008647213.3} & 2,386 & 52.22\% / 41.38\% & 0.0921 & 0.20855 & \texttt{N28 (76/76)} \\
\bottomrule
\end{longtable}

\begin{thebibliography}{99}
\bibitem{ref1} Ponce-Toledo, R. I., Deschamps, P., López-García, P., Zivanovic, Y., Benzerara, K., \& Moreira, D. (2017). "An early-branching freshwater cyanobacterium at the origin of plastids." Current Biology, 27(3), 386-391. \url{https://doi.org/10.1016/j.cub.2016.11.056}
\bibitem{ref2} Sánchez-Baracaldo, P., Raven, J. A., Pisani, D., \& Knoll, A. H. (2017). "Early photosynthetic eukaryotes inhabited low-salinity habitats." Proceedings of the National Academy of Sciences, 114(37), E7737-E7745. \url{https://doi.org/10.1073/pnas.1620089114}
\bibitem{ref3} Gibson, T. M., Shih, P. M., Cumming, V. M., Fischer, W. W., Crockford, P. W., Hodgskiss, M. S., ... \& Halverson, G. P. (2018). "Precise age of Bangiomorpha pubescens dates the origin of eukaryotic photosynthesis." Geology, 46(2), 135-138. \url{https://doi.org/10.1130/G39829.1}
\bibitem{ref4} Bengtson, S., Sallstedt, T., Belivanova, V., \& Whitehouse, M. (2017). "Three-dimensional preservation of cellular and subcellular structures suggests 1.6 billion-year-old crown-group red algae." PLOS Biology, 15(3), e2000735. \url{https://doi.org/10.1371/journal.pbio.2000735}
\bibitem{ref5} Miao, L., Yin, Z., Knoll, A. H., Qu, Y., \& Zhu, M. (2024). "1.63-billion-year-old multicellular eukaryotes from the Chuanlinggou Formation in North China." Science Advances, 10(4), eadk3208. \url{https://doi.org/10.1126/sciadv.adk3208}
\bibitem{ref6} Tang, Q., Pang, K., Yuan, X., \& Xiao, S. (2020). "A one-billion-year-old multicellular chlorophyte." Nature Ecology \& Evolution, 4(4), 543-549. \url{https://doi.org/10.1038/s41559-020-1122-9}
\bibitem{ref7} Nowack, E. C., Melkonian, M., \& Glöckner, G. (2008). "Chromatophore genome sequence of Paulinella sheds light on acquisition of photosynthesis by eukaryotes." Current Biology, 18(6), 410-418. \url{https://doi.org/10.1016/j.cub.2008.02.051}
\bibitem{ref8} Lhee, D., Ha, J. S., Kim, S., Park, M. G., Bhattacharya, D., \& Yoon, H. S. (2019). "Evolutionary dynamics of the chromatophore genome in three photosynthetic Paulinella species." Scientific Reports, 9(1), 2560. \url{https://doi.org/10.1038/s41598-019-38621-8}
\bibitem{ref9} Liu, H., Li, A., Rochaix, J. D., \& Liu, Z. (2023). "Architecture of chloroplast TOC-TIC translocon supercomplex." Nature, 615(7951), 349-357. \url{https://doi.org/10.1038/s41586-023-05744-y}
\bibitem{ref10} Jin, Z., Wan, L., Zhang, Y., Li, X., Cao, Y., Liu, H., ... \& Yan, Z. (2022). "Structure of a TOC-TIC supercomplex spanning two chloroplast envelope membranes." Cell, 185(25), 4788-4800. \url{https://doi.org/10.1016/j.cell.2022.10.030}
\bibitem{ref11} Kikuchi, S., Bédard, J., Hirano, M., Hirabayashi, Y., Oishi, M., Imai, M., ... \& Nakai, M. (2013). "Uncovering the protein translocon at the chloroplast inner envelope membrane." Science, 339(6119), 571-574. \url{https://doi.org/10.1126/science.1229262}
\bibitem{ref12} Day, P. M., Potter, L. S., Inoue, K., \& Theg, S. M. (2014). "Chloroplast outer membrane beta-barrel protein Toc75 requires POTRA domains for protein import." Plant Cell, 26(7), 2997-3009. \url{https://doi.org/10.1105/tpc.114.126896}
\bibitem{ref13} Jackson, C., Clayden, S., \& Reyes-Prieto, A. (2015). "The Glaucophyta: the blue-green plants in a nutshell." Acta Societatis Botanicorum Poloniae, 84(2), 149-165. \url{https://doi.org/10.5586/asbp.2015.020}
\bibitem{ref14} Price, D. C., Chan, C. X., Yoon, H. S., Yang, E. C., Qiu, H., Weber, A. P., ... \& Bhattacharya, D. (2012). "Cyanophora paradoxa genome elucidates origin of photosynthesis in algae and plants." Science, 335(6070), 843-847. \url{https://doi.org/10.1126/science.1213561}
\bibitem{ref15} Morris, J. L., Puttick, M. N., Clark, J. W., Edwards, D., Kenrick, P., Pressel, S., ... \& Donoghue, P. C. (2018). "The timescale of early land plant evolution." Proceedings of the National Academy of Sciences, 115(10), E2274-E2283. \url{https://doi.org/10.1073/pnas.1719588115}
\bibitem{ref16} Del Cortona, A., Jackson, C. J., Bucchini, F., Van Bel, M., D’hondt, S., Škaloud, P., ... \& De Clerck, O. (2020). "Neoproterozoic origin and multiple transitions to macroscopic growth in green seaweeds." Proceedings of the National Academy of Sciences, 117(5), 2551-2559. \url{https://doi.org/10.1073/pnas.1910060117}
\bibitem{ref17} Yoon, H. S., Hackett, J. D., Ciniglia, C., Pinto, G., \& Bhattacharya, D. (2004). "A molecular timeline for the origin of photosynthetic eukaryotes." Molecular Biology and Evolution, 21(5), 809-818. \url{https://doi.org/10.1093/molbev/msh075}
\bibitem{ref18} Bekker, A., Holland, H. D., Wang, P. L., Rumble, D., Stein, H. J., Hannah, J. L., ... \& Beukes, N. J. (2004). "Dating the rise of atmospheric oxygen." Nature, 427(6970), 117-120. \url{https://doi.org/10.1038/nature02260}
\bibitem{ref19} Hodgskiss, M. S., Crockford, P. W., Peng, Y., Wing, B. A., \& Horner, T. J. (2019). "A productivity collapse to end Earth’s Great Oxidation." Proceedings of the National Academy of Sciences, 116(35), 17207-17212. \url{https://doi.org/10.1073/pnas.1900325116}
\bibitem{ref20} Tomitani, A., Knoll, A. H., Cavanaugh, C. M., \& Ohno, T. (2006). "The evolutionary diversification of cyanobacteria: molecular-phylogenetic and paleontological perspectives." Proceedings of the National Academy of Sciences, 103(14), 5442-5447. \url{https://doi.org/10.1073/pnas.0600999103}
\bibitem{ref21} Nakamura, Y., Kaneko, T., Sato, S., Mimuro, M., Miyashita, H., Tsuchiya, T., ... \& Tabata, S. (2003). "Complete genome structure of Gloeobacter violaceus PCC 7421, a cyanobacterium that lacks thylakoids." DNA Research, 10(4), 137-145. \url{https://doi.org/10.1093/dnares/10.4.137}
\bibitem{ref22} Saw, J. H., Schatz, M., Brown, M. V., Kunkel, D. D., Foster, J. S., Shick, H., ... \& Donachie, S. P. (2013). "Cultivation and complete genome sequencing of Gloeobacter kilaueensis sp. nov., from a lava cave in Kilauea Caldera, Hawaii." PLOS ONE, 8(10), e76376. \url{https://doi.org/10.1371/journal.pone.0076376}
\bibitem{ref23} Couradeau, E., Benzerara, K., Gérard, E., Moreira, D., Bernard, S., Brown, G. E., \& López-García, P. (2012). "An early-branching microbialite cyanobacterium forms intracellular carbonates." Science, 336(6080), 459-462. \url{https://doi.org/10.1126/science.1218178}
\bibitem{ref24} Osteryoung, K. W., \& Pyke, K. A. (2014). "Division and dynamic morphology of plastids." Annual Review of Plant Biology, 65, 443-472. \url{https://doi.org/10.1146/annurev-arplant-050213-035748}
\bibitem{ref25} Andersson, I., \& Backlund, A. (2008). "Structure and function of Rubisco." Plant Physiology and Biochemistry, 46(3), 275-291. \url{https://doi.org/10.1016/j.plaphy.2008.01.001}
\bibitem{ref26} Katoh, K., \& Standley, D. M. (2013). "MAFFT multiple sequence alignment software version 7: improvements in performance and usability." Molecular Biology and Evolution, 30(4), 772-780. \url{https://doi.org/10.1093/molbev/mst010}
\bibitem{ref27} Minh, B. Q., Schmidt, H. A., Chernomor, O., Schrempf, D., Woodhams, M. D., Von Haeseler, A., \& Lanfear, R. (2020). "IQ-TREE 2: new models and efficient methods for phylogenetic inference in the genomic era." Molecular Biology and Evolution, 37(5), 1530-1534. \url{https://doi.org/10.1093/molbev/msaa015}
\bibitem{ref28} To, T. H., Jung, M., Lycett, S., \& Gascuel, O. (2016). "Fast dating using least-squares criteria and algorithms." Systematic Biology, 65(1), 82-97. \url{https://doi.org/10.1093/sysbio/syv068}
\bibitem{ref29} Yang, Z. (2007). "PAML 4: phylogenetic analysis by maximum likelihood." Molecular Biology and Evolution, 24(8), 1586-1591. \url{https://doi.org/10.1093/molbev/msm088}
\bibitem{ref30} Goldman, N., \& Yang, Z. (1994). "A codon-based model of nucleotide substitution for protein-coding DNA sequences." Molecular Biology and Evolution, 11(5), 725-736. \url{https://doi.org/10.1093/oxfordjournals.molbev.a040153}
\end{thebibliography}

\end{document}
